\chapter{Fundamente teoretice și contextul actual}
\label{chap:ch2}

\par Prezentul capitol stabilește fundamentul teoretic necesar înțelegerii și proiectării arhitecturii sistemului Sentinel. Sunt analizate mecanismele intrinseci de securitate ale nucleului Linux, inovațiile aduse de tehnologia eBPF, taxonomia sistemelor de detectare a intruziunilor, aparatul matematic al algoritmului Isolation Forest, precum și o analiză comparativă a soluțiilor consacrate în domeniu.


\section{Modelul de securitate Linux}
\label{sec:linux_security}

\subsection{Apeluri de sistem și frontiera kernel-utilizator}
\label{subsec:syscalls}

\par Granița fundamentală de izolare între codul cu privilegii reduse (spațiul utilizator / user-space) și codul cu privilegii absolute (nucleul / kernel-space) în sistemele Linux este definită de interfața apelurilor de sistem (syscalls). Orice acțiune care necesită interacțiunea cu resursele hardware sau logice ale sistemului --- instanțierea de conexiuni (\texttt{socket}), crearea de procese, manipularea sistemului de fișiere, alterarea privilegiilor --- trebuie să traverseze această interfață printr-o tranziție de context strict controlată \cite{Kerrisk2010}.

\par Din perspectiva securității cibernetice, monitorizarea apelurilor de sistem oferă o vizibilitate exhaustivă și imposibil de eludat asupra comportamentului real al unui proces: o rutină malițioasă care încearcă să încarce un modul de kernel va invoca inevitabil \texttt{finit\_module} (syscall 313), o tentativă de escaladare a privilegiilor va apela \texttt{setuid} (syscall 105) sau \texttt{capset} (syscall 126), în timp ce exfiltrarea datelor va fi tradusă prin generarea de apeluri \texttt{socket} și \texttt{connect} (syscall-urile 41, respectiv 42) \cite{Love2010}.

\par Kernel-ul Linux x86\_64 expune peste 300 de apeluri de sistem, dar numai un subansamblu specific prezintă interes direct din perspectiva securității. Sistemul Sentinel monitorizează 20 de apeluri de sistem, selectate pe baza incidenței lor empirice în anatomia malware-ului pentru Linux \cite{Cozzi2018}, acoperind categoriile: execuție de procese, modificare de privilegii, operații de memorie, activitate de rețea, operații pe fișiere sensibile și management de module de kernel. Tabelul~\ref{tab:syscalls_risk} prezintă o selecție cu ponderile de risc corespunzătoare.

\begin{table}[htbp]
\begin{center}
\begin{tabular}{|p{80pt}|p{40pt}|p{90pt}|p{60pt}|}
\hline
\textbf{Syscall} & \textbf{Nr.} & \textbf{Categorie alertă} & \textbf{Pondere} \\
\hline\hline
\texttt{finit\_module} & 313 & HiddenPID           & 95 \\
\hline
\texttt{init\_module}  & 175 & HiddenPID           & 95 \\
\hline
\texttt{delete\_module}& 176 & HiddenPID           & 85 \\
\hline
\texttt{ptrace}        & 101 & ProcessInjection    & 75 \\
\hline
\texttt{setuid}        & 105 & PrivilegeEscalation & 70 \\
\hline
\texttt{setgid}        & 106 & PrivilegeEscalation & 70 \\
\hline
\texttt{capset}        & 126 & PrivilegeEscalation & 70 \\
\hline
\texttt{execve}        & 59  & ParentChildAnomaly  & 50 \\
\hline
\texttt{mprotect}      & 10  & MemoryInjection     & 45 \\
\hline
\texttt{mmap}          & 9   & MemoryInjection     & 40 \\
\hline
\texttt{connect}       & 42  & NetworkAnomaly      & 15 \\
\hline
\texttt{socket}        & 41  & NetworkAnomaly      & 10 \\
\hline
\end{tabular}
\end{center}
\caption{Selecție de apeluri de sistem monitorizate de Sentinel cu ponderile de risc asociate}
\label{tab:syscalls_risk}
\end{table}


\subsection{Linux Security Modules}
\label{subsec:lsm}

\par Cadrul Linux Security Modules (LSM), teoretizat și implementat de Wright et al. \cite{Wright2002}, pune la dispoziție o arhitectură bazată pe cârlige de interceptare (hooks) plasate strategic în punctele decizionale ale nucleului. Această infrastructură permite modulelor de securitate (ex. SELinux, AppArmor) să impună politici stricte de tip Mandatory Access Control (MAC).

\par O limitare inerentă a modulelor LSM tradiționale este natura lor declarativă și statică, nefiind concepute pentru analiza comportamentală dinamică a secvențelor de execuție. Din acest motiv, un sistem de detectare a anomaliilor comportamentale, precum Sentinel, reprezintă o completare necesară și ortogonală față de mecanismele LSM. Deși eBPF permite atașarea programelor la punctele de interceptare LSM (prin \texttt{BPF\_PROG\_TYPE\_LSM}, funcționalitate introdusă în Linux 5.7), Sentinel utilizează tehnologia kprobes pentru a asigura o compatibilitate retroactivă mai largă pe nucleele $\geq 5.4$.


\subsection{Namespaces și grupuri de control (cgroups)}
\label{subsec:cgroups}

\par Mecanismele interne ale nucleului (PID, net, mount, user, UTS, IPC) asigură izolarea stării globale a kernel-ului în contexte virtualizate pentru grupuri distincte de procese, reprezentând fundamentul tehnologiilor moderne de containerizare \cite{Kerrisk2010}. În paralel, mecanismul Control Groups v2 (cgroup v2), fundamentat inițial de Menage \cite{Menage2007} și standardizat în nucleul 4.5, gestionează alocarea, limitarea și monitorizarea resurselor fizice (CPU, memorie, I/O bloc) pentru ierarhii de procese.

\par Spre deosebire de iterația anterioară (cgroup v1), arhitectura v2 propune o ierarhie unificată și un model de delegare mai coerent, eliminând ambiguitățile de configurare care apăreau la apartenența simultană a unui proces la sub-arbori de control diferiți. Arhitectura de izolare a sistemului Sentinel impune la pornire prezența controllerelor \texttt{cpu} și \texttt{memory} în \texttt{/sys/fs/cgroup/cgroup.controllers}, condiție satisfăcută nativ pe distribuțiile Debian 12 și Ubuntu 22.04+.


\section{Extended Berkeley Packet Filter (eBPF)}
\label{sec:ebpf}

\subsection{Evoluție de la BPF clasic la eBPF}
\label{subsec:ebpf_history}

\par Berkeley Packet Filter (BPF) a fost conceptualizat de McCanne și Jacobson în 1993, vizând dezvoltarea unei arhitecturi extrem de eficiente pentru capturarea și filtrarea traficului de rețea, operând printr-o mașină virtuală minimalistă pe 32 de biți \cite{McCanne1993}. Schimbarea de paradigmă pe care a propus-o BPF a fost executarea logicii de filtrare direct în kernel, evaluând pachetele la sursă pentru a preveni transferul inutil de date peste frontiera privilegiată.

\par Începând cu Linux 3.18 (2014), această arhitectură a suferit o extensie masivă sub acronimul eBPF: lărgirea registrelor la 64 de biți, îmbogățirea setului de instrucțiuni, introducerea structurilor de date partajate (maps) și capacitatea de atașare asincronă la aproape orice funcție a sistemului de operare \cite{Rice2023}. O inovație critică o reprezintă Verificatorul Formal (The Verifier), un mecanism care analizează codul înainte de execuție pentru a garanta finalizarea acestuia (lipsa buclelor infinite) și respectarea strictă a limitelor de memorie accesată. Compilatorul JIT (Just-In-Time) convertește ulterior bytecode-ul validat în cod mașină nativ, reducând la minimum costul de performanță (overhead).

\par Ecosistemul eBPF a evoluat rapid: bibliotecile \texttt{libbpf} și \texttt{cilium/ebpf} \cite{CiliumEBPF} oferă o interfață Go/C de nivel înalt pentru încărcarea și gestionarea programelor eBPF, iar CO-RE (Compile Once, Run Everywhere) rezolvă problema portabilității între versiuni de kernel prin relocări BTF.


\subsection{Maps eBPF și mecanismul ring buffer}
\label{subsec:ebpf_maps}

\par Maps eBPF sunt structuri de stocare persistente, expuse asincron atât spațiului kernel, cât și celui utilizator. Arhitectura Sentinel implementează trei structuri fundamentale:

\begin{itemize}
    \item \texttt{BPF\_MAP\_TYPE\_HASH}: tabelă de dispersie alocată pentru păstrarea excluderilor (\texttt{pid\_blacklist}).
    \item \texttt{BPF\_MAP\_TYPE\_PERCPU\_ARRAY}: vector instanțiat per-CPU pentru o contorizare atomică, fără blocaje (lock-free), a evenimentelor suprimate (\texttt{drop\_count}).
    \item \texttt{BPF\_MAP\_TYPE\_RINGBUF}: coadă circulară de mare viteză, dimensionată la 64 MiB, utilizată ca bus principal de transmitere a evenimentelor interceptate.
\end{itemize}

\par Ring buffer-ul (introdus în kernel 5.8) oferă avantaje față de perf buffer-ul precedent: elimină o copiere suplimentară de date, suportă rezervarea atomică a spațiului urmată de scriere directă și notificarea consumatorului printr-un descriptor de fișier abilitat în acest sens prin mecanismul epoll. Gregg \cite{Gregg2019} documentează că ring buffer-ul reduce latența de notificare cu 20--40\% față de perf buffer în scenariile de evenimente cu frecvență ridicată.


\subsection{Kprobes: instrumentarea dinamică a kernel-ului}
\label{subsec:kprobes}

\par Prin intermediul mecanismului kprobes (Kernel Probes), putem instrumenta codul nucleului în timp real, injectând rutine de tratare exact acolo unde avem nevoie, fără a atinge codul sursă și fără a recompila sistemul [CRKH05]. Funcționarea sa este elegantă: sonda înlocuiește discret instrucțiunea originală cu o întrerupere software (de tipul INT3 pe x86). Când execuția lovește această „capcană”, controlul este pasat instantaneu mașinii virtuale eBPF. După procesarea evenimentului, sistemul restaurează și execută instrucțiunea originală. Deși pare un proces complex, această deviere este extrem de rapidă, inducând o întârziere de doar 100–300 nanosecunde [VCP+20] – un cost insesizabil în contextul unui apel de sistem obișnuit.

\par Sentinelul atașează kprobe-uri pe funcțiile \texttt{\_\_x64\_sys\_\{name\}} corespunzătoare celor 20 de syscall-uri monitorizate. Atașarea este dinamică: la pornire, daemonul iterează peste toate programele din colecția eBPF cu prefix \texttt{kp\_} și le atașează la funcțiile kernel corespunzătoare, fără a fi necesară enumerarea lor explicită în codul Go.

\section{Sisteme de detectare a intruziunilor}
\label{sec:ids}

\subsection{Taxonomie și clasificare}
\label{subsec:ids_taxonomy}

\par Conceptul monitorizării securității a fost formalizat inițial de Anderson (1980) ca o metodă de prelucrare a jurnalelor de audit în scopul detectării utilizării neautorizate a sistemelor informatice \cite{Anderson1980}. Ulterior, Denning (1987) a publicat fundamentul teoretic al sistemelor IDS (sisteme de detecție a intruziunilor), propunând un model independent de sistem capabil să evalueze statistic deviațiile de la comportamentul istoric normal \cite{Denning1987}.

\par Taxonomia propusă de Axelsson (2000) \cite{Axelsson2000} oferă un cadru de clasificare în care Sentinel se raportează astfel:

\begin{itemize}
    \item \textit{Sursa datelor}: HIDS (Host-based) versus NIDS (Network-based). Sentinel este exclusiv un sistem HIDS, rezoluția sa fiind limitată la apelurile de sistem ale mașinii gazdă.
    \item \textit{Metoda de detecție}: Detecție bazată pe semnături versus anomalii. Sistemul de față abordează problema hibrid, combinând un set de euristici statice cu un model nesupervizat pentru anomalii.
    \item \textit{Frecvența analizei}: detecție continuă versus periodică. Sentinel operează în timp real.
    \item \textit{Tipul de răspuns}: alertare pasivă versus limitare activă (containment). Sentinel oferă un răspuns activ, autonom și gradual.
\end{itemize}


\subsection{Detecție bazată pe anomalii versus semnături}
\label{subsec:anomaly_vs_signature}

\par Abordarea bazată pe semnături excelează prin rata infimă a rezultatelor fals pozitive pentru atacurile cunoscute, însă suferă din cauza ineficienței totale în fața atacurilor neclasificate (zero-day) \cite{GarciaTe2009}. Alternativa, detecția anomaliilor, profilează starea de normalitate a sistemului; deși superioară în identificarea vectorilor de atac noi, aceasta este penalizată de un volum mai mare de rezultate fals pozitive \cite{Chandola2009}.

\begin{comment}
TODO: check factuality further
\par În stabilirea algoritmului adecvat, studiul lui Goldstein și Uchida \cite{Goldstein2016} demonstrează performanța net superioară a modelului Isolation Forest pe distribuții de date unde anomaliile reprezintă un procent extrem de redus ($< 5\%$), modelând perfect realitatea evenimentelor de securitate. Această alegere este validată și de lucrarea de sinteză elaborată de Buczak și Guven \cite{Buczak2016}, care evidențiază eficiența metodelor de ansamblu (ensemble methods) pe spații de date cu dimensionalitate variabilă.
\end{comment}

\section{Algoritmul Isolation Forest}
\label{sec:isolation_forest}

\par Isolation Forest (iForest), formulat de Liu, Ting și Zhou \cite{Liu2008IF, Liu2012IF}, schimbă paradigma clasică a detecției anomaliilor: în loc să profileze marea masă a datelor normale, algoritmul încearcă să izoleze explicit valorile aberante. Postulatul principal este că anomaliile reprezintă instanțe rare și divergente structural; prin urmare, o serie de partiționări aleatorii (random splits) va izola un punct anormal mult mai aproape de rădăcina unui arbore de decizie comparativ cu un punct benign.

\par Formal, dat fiind un set de antrenament $X = \{x_1, \ldots, x_n\} \subset \mathbb{R}^d$, construirea unui arbore de izolare presupune selecția aleatoare a unui atribut $q \in \{1, \ldots, d\}$ și partiționarea spațiului pe un punct $p$ ales dintr-o distribuție uniformă în intervalul $[\min_q(X), \max_q(X)]$. Evaluarea (scorul de anomalie) pentru un punct $x$ este obținută prin relația:

\begin{equation}
    s(x, n) = 2^{-\frac{E[h(x)]}{c(n)}}
    \label{eq:isolation_score}
\end{equation}

\noindent unde $E[h(x)]$ reprezintă valoarea așteptată (adâncimea medie) a lungimii căii către punctul $x$ generată pe ansamblul de arbori, iar $c(n) = 2H(n-1) - \frac{2(n-1)}{n}$ este adâncimea medie asimptotică într-un arbore binar de căutare nereușită, cu $H$ funcția armonică \cite{Liu2012IF}. Scorul variază în intervalul $(0, 1)$: valori aproape de $1$ indică anomalii puternice, valori aproape de $0.5$ indică puncte normale sau ambigue.

\par Deoarece scorul brut emis de model (valoarea funcției de decizie \texttt{decision\_function}, negativă pentru anomalii) necesită o conversie pe scala operațională $[0, 100]$, Sentinel aplică o transformare sigmoidă:

\begin{equation}
    \text{threat\_score} = \frac{100}{1 + e^{-25 \cdot (-\text{raw\_score})}}
    \label{eq:sigmoid_transform}
\end{equation}

\par Factorul de scalare 25 a fost ales empiric pentru a produce o separare clară între punctele normale (\texttt{threat\_score} $< 30$) și anomalii moderate (\texttt{threat\_score} $\in [50, 70]$), menținând în același timp sensibilitate pentru valorile extreme.

\par Complexitatea antrenării este $\mathcal{O}(t \cdot \psi \cdot \log \psi)$, cu $t$ numărul de arbori și $\psi$ dimensiunea sub-eșantionării (subsampling), în timp ce faza de inferență se execută în timp $\mathcal{O}(t \cdot \log \psi)$ per eveniment --- acceptabilă pentru bugetul de latență impus de constrângerile decizionale în timp real (RN2).


\section{Instrumente existente și lucrări corelate}
\label{sec:related_work}

\par \textbf{Falco} \cite{FalcoProject}, proiect incubat de CNCF și derivat din Sysdig, este cel mai apropiat echivalent open-source al Sentinelului: utilizează eBPF sau un driver de kernel pentru monitorizarea syscall-urilor și permite definirea de reguli de alertă în YAML. Spre deosebire de Sentinel, Falco nu implementează un scorer ML integrat, nu aplică acțiuni automate de containment (răspunsul este delegat unor sisteme externe precum Falcosidekick) și nu menține o bază de date locală a evenimentelor pentru analiza retrospectivă.

\par \textbf{auditd} \cite{LinuxAudit} este componenta de audit a kernel-ului Linux, standardizată și prezentă în toate distribuțiile majore. Jurnalizează apeluri de sistem pe baza unor reguli configurate static. Din pricina prelucrării asincrone și a interacțiunii constante cu sistemul de fișiere pentru stocarea jurnalelor, auditd induce o latență incompatibilă cu acțiunile de izolare în timp real, limitându-se la rolul de analiză retrospectivă (forensics).

\par \textbf{Virtual Machine Introspection (VMI)}: Metodologia teoretizată de Garfinkel și Rosenblum (2003) \cite{Garfinkel2003} propune securizarea prin mutarea monitorului în hipervizor, complet izolat de sistemul de operare monitorizat. Implementări precum DRAKVUF \cite{Lengyel2014} utilizează extensiile de virtualizare (Xen) pentru a analiza memoria mașinii virtuale din exterior. Limita severă a acestei abordări este dependența de virtualizarea hardware, fiind incompatibilă cu mediile bare-metal sau infrastructurile cloud moderne bazate pe containere.

\par \textbf{kGuard} \cite{Kemerlis2012} abordează o problemă complementară --- protecția kernel-ului împotriva atacurilor de tip return-to-user --- prin instrumentarea statică a codului kernel la compilare. Aceste abordări de securitate prin compilare sunt ortogonale față de Sentinel, care se concentrează pe detectarea comportamentului la runtime.

\par O constatare empirică importantă din literatura de specialitate provine din Cozzi et al. \cite{Cozzi2018}: 70\% din eșantioanele de malware Linux analizate utilizează socket-uri de rețea, 30\% apelează \texttt{ptrace} sau mapează pagini cu drepturi de execuție, iar 18\% utilizează module de kernel. Aceste statistici au constituit criteriul de selecție și de ponderare a syscall-urilor monitorizate de Sentinel.


\section{Peisajul amenințărilor în ecosistemul Linux și Cadrul MITRE ATT\&CK}
\label{sec:threat_landscape}

\par Evaluarea oricărui sistem modern de detectare a intruziunilor necesită raportarea riguroasă la un cadru standardizat de amenințări. În ecosistemul Linux, vectorii de atac au evoluat de la simple scripturi de exploatare la arhitecturi complexe, mapate exhaustiv în cadrul MITRE ATT\&CK \cite{MITRE_ATTCK}. Evaluarea eficacității sistemului Sentinel se fundamentează pe o analiză teoretică a 20 de scenarii de atac care acoperă tactici critice: persistență, escaladarea privilegiilor, evaziune defensivă, acces la credențiale și execuție.

\par În continuare sunt detaliate principalele categorii teoretice de amenințări care justifică arhitectura de monitorizare a sistemului:

\subsection{Rootkit-uri și persistența la nivel de nucleu}
\label{subsec:rootkits}

\par Compromiterea nucleului Linux s-a realizat istoric prin rootkit-uri de tip LKM (\textit{Loadable Kernel Module}), care vizează alterarea fluxului de execuție din spațiul privilegiat. Un etalon istoric este rootkit-ul \textit{Diamorphine}, apărut în 2013, care interceptează apeluri de sistem (precum \texttt{getdents64}) pentru a ascunde procese și asigură o portiță de acces (\textit{backdoor}) gestionată prin abuzarea apelului de sistem kill și interceptarea unor semnale POSIX predefinite pentru a declanșa funcții ascunse, care ar putea da privilegii atacatorului. Încărcarea acestor module este semnalizată primar de apelul \texttt{init\_module}.

\par Pentru a evita lăsarea de artefacte pe disc, atacatorii utilizează apelul \texttt{finit\_module}, care permite încărcarea unui modul compilat direct dintr-un descriptor de fișier menținut exclusiv în memoria volatilă. Amenințările de generație nouă includ rootkit-urile eBPF, exemplificate de modelul teoretizat de \textit{TripleCross} în 2022 \cite{TripleCross2022}. Acestea demonstrează capacitatea atacatorilor de a subverti însăși tehnologia eBPF pentru a intercepta apeluri de sistem și a orchestra injecții de cod invizibile scanărilor tradiționale.

\subsection{Escaladarea privilegiilor}
\label{subsec:privesc}

\par Obținerea drepturilor absolute de administrator (UID 0) pornind de la un cont restricționat reprezintă un obiectiv primar în faza de post-exploatare. Pe sistemele Linux, această acțiune este mediată în principal de apelurile \texttt{setuid} (pentru alterarea directă a identificatorului de utilizator) și \texttt{capset} (pentru modificarea fină a capabilităților POSIX). Un exemplu elocvent al severității acestei clase de atacuri este vulnerabilitatea \textit{PwnKit} (CVE-2021-4034), care a exploatat o corupție de memorie în componenta PolicyKit pentru a permite execuția cu succes a funcției \texttt{setuid(0)} pe majoritatea distribuțiilor majore \cite{PwnKit2022}.

\subsection{Evaziunea defensivă și execuția \textit{fileless}}
\label{subsec:fileless}

\par Malware-ul de tip \textit{fileless} execută instrucțiuni exclusiv în memorie, ocolind soluțiile antivirus tradiționale fundamentate pe analiza sistemului de fișiere. Această tehnică a fost facilitată structural de introducerea apelului \texttt{memfd\_create} în nucleul Linux 3.17 (2014), care permite crearea unui fișier anonim susținut doar de memoria RAM. Atacatorul poate scrie codul malițios (de exemplu, un binar ELF) în acest descriptor și îl poate executa prin apelul \texttt{execve} (referențiind calea virtuală \texttt{/proc/self/fd/}), fără ca executabilul să persiste vreodată pe un mediu de stocare fizic \cite{Gruss2016}.

\subsection{Injectarea de procese}
\label{subsec:process_injection}

\par Apelul de sistem \texttt{ptrace}, conceput inițial ca mecanism legitim pentru depanare (\textit{debugging}), a devenit instrumentul canonic pentru tehnicile de injectare a proceselor. Acesta permite unui proces malițios să suspende o țintă legitimă (\texttt{PTRACE\_ATTACH}), să îi altereze pointerul de instrucțiune (\texttt{PTRACE\_SETREGS}) și să scrie \textit{shellcode} direct în spațiul său de adresare (\texttt{PTRACE\_POKETEXT}).

\par O derivare complexă a acestei tehnici este golirea proceselor (\textit{Process Hollowing} sau \textit{RunPE}). Un proces legitim este instanțiat și suspendat, i se alocă o regiune de memorie nouă, simultan lizibilă, scriptibilă și executabilă (RWX) prin \texttt{mmap}, iar fluxul normal este înlocuit cu logica atacatorului \cite{Elastic2022}.

\subsection{Accesul la credențiale și controlul (command-and-control / c2)}
\label{subsec:c2}

\par \textbf{Furtul de credențiale.} Atacatorii enumeră sistematic fișiere sensibile (de exemplu, \texttt{/etc/shadow}, chei SSH private) printr-o cascadă de apeluri \texttt{openat}. Pentru extragerea parolelor aflate în uz, tehnici similare utilitarului Mimikatz (Windows) sunt aplicate pe Linux prin citirea memoriei proceselor rulante, mapată în \texttt{/proc/<pid>/mem} \cite{Bundles2021}.

\par \textbf{Exfiltrarea datelor.} Stabilirea canalelor de comandă și control implică adesea instanțierea de \textit{reverse shells}, procedeu în care gazda compromisă inițiază conexiuni TCP de ieșire (\textit{outbound}) către atacator, ocolind astfel regulile restrictive de firewall pentru traficul de intrare. Amprenta comportamentală constă în serii de apeluri \texttt{socket} și \texttt{connect}, urmate de deturnarea fluxurilor de intrare/ieșire prin invocarea unui shell via \texttt{execve} \cite{Cozzi2018}.